%----------------------------------------------------------------------------------------------------%
% Dokumentspezifische Befehle (nur für dieses Wörterbuch, deshalb lokal und nicht in newcommands.tex)
%----------------------------------------------------------------------------------------------------%

% \eintrag{Marke}{Artikel}{Begriff}{english term}{Lautschrift EN}{Zeichen/Skizze}
% Schaltet das babel-Kürzel " aus, BEVOR die Argumente gelesen werden – dadurch darf die Lautschrift
% wie bei \ipa{...} mit " (Hauptbetonung) und "" (Nebenbetonung) geschrieben werden.
\makeatletter
\newcommand\eintrag{\begingroup\shorthandoff{"}\eintrag@aux}
\newcommand\eintrag@aux[6]{\endgroup
	#2 \textbf{#3}\phantomsection\label{e:#1}
	& \textbf{#4}\newline{\small\color{black!70}\ipa{#5}}
	& #6 \\ \midrule}
\makeatother

% Tabellenkopf für jedes Kapitel (wird auf jeder Folgeseite wiederholt)
\newenvironment{woerter}{%
	\arrayrulecolor{black!20}%
	\begin{longtable}{>{\raggedright\arraybackslash}m{5.3cm} >{\raggedright\arraybackslash}m{6.1cm} >{\centering\arraybackslash}m{4.6cm}}
		\rowcolor{cornflower!15}\textbf{Deutsch} & \textbf{Englisch} & \textbf{Zeichen \,/\, Skizze} \\ \midrule
		\endhead}{%
	\end{longtable}\arrayrulecolor{black}}

% Kleine Skizzen (Gestaltung wie im Skill latex-geometrie: Figur blau very thick, Maße/Hervorhebung rot)
\tikzset{
	fig/.style={blue, very thick},
	mass/.style={red, thick},
	hilf/.style={dashed},
	achse/.style={-{Stealth[length=4pt]}, thin},
}
\newcommand{\skizze}[1]{\tikz[baseline=(current bounding box.center), scale=0.45, line cap=round,
	line join=round, every node/.style={font=\scriptsize}, outer sep=0pt]{#1}}

% Grundbegriffe
\newcommand{\skPunkt}{\skizze{\fill[blue] (0,0) circle (2.5pt); \node[above right] at (0,0) {$P$};}}
\newcommand{\skGerade}{\skizze{\draw[fig] (0,0) -- (3.6,0.8); \node[below right] at (3.2,0.7) {$g$};}}
\newcommand{\skStrecke}{\skizze{\draw[fig] (0,0) -- (3.2,0.7); \fill (0,0) circle (2pt) (3.2,0.7) circle (2pt);
	\node[below] at (0,0) {$A$}; \node[below] at (3.2,0.7) {$B$};}}
\newcommand{\skStrahl}{\skizze{\draw[fig, -{Stealth}] (0,0) -- (3.6,0.8); \fill (0,0) circle (2pt);
	\node[below] at (0,0) {$S$};}}
\newcommand{\skParallel}{\skizze{\draw[fig] (0,0) -- (3.5,0.6) node[right] {$g$};
	\draw[fig] (0,-0.9) -- (3.5,-0.3) node[right] {$h$};}}
\newcommand{\skSenkrecht}{\skizze{\draw[fig] (-1.6,0) -- (1.6,0) node[right] {$g$};
	\draw[fig] (0,-0.5) -- (0,1.8) node[right] {$h$};
	\draw[mass] (0.45,0) arc[start angle=0, end angle=90, radius=0.45]; \fill[red] (45:0.24) circle (1pt);}}
\newcommand{\skAbstand}{\skizze{\draw[fig] (-0.4,0) -- (3.2,0) node[right] {$g$}; \fill (1.4,1.6) circle (2pt) node[above] {$P$};
	\draw[mass, hilf] (1.4,1.6) -- node[right] {$d$} (1.4,0);
	\draw[mass] (1.75,0) arc[start angle=0, end angle=90, radius=0.35]; \fill[red] (1.4,0) ++(45:0.19) circle (0.8pt);}}

% Winkel
\newcommand{\skWinkel}[1]{\skizze{\draw[mass] (0.7,0) arc[start angle=0, end angle=#1, radius=0.7];
	\draw[fig] (2.3,0) -- (0,0) -- (#1:2.3); \node at ({#1/2}:1.15) {$\alpha$};}}
\newcommand{\skRechterWinkel}{\skizze{\draw[mass] (0.6,0) arc[start angle=0, end angle=90, radius=0.6];
	\fill[red] (45:0.32) circle (1pt); \draw[fig] (2.3,0) -- (0,0) -- (0,2.3);}}
\newcommand{\skVollwinkel}{\skizze{\draw[mass] (0.6,0) arc[start angle=0, end angle=355, radius=0.6];
	\draw[fig] (0,0) -- (2.3,0); \fill[blue] (0,0) circle (1.5pt); \node at (135:1.0) {$\alpha$};}}
\newcommand{\skNebenwinkel}{\skizze{\draw[mass] (0.7,0) arc[start angle=0, end angle=60, radius=0.7];
	\draw[mass] (60:0.55) arc[start angle=60, end angle=180, radius=0.55];
	\draw[fig] (-2.2,0) -- (2.2,0); \draw[fig] (0,0) -- (60:2.2);
	\node at (28:1.15) {$\alpha$}; \node at (125:1.0) {$\beta$};}}
\newcommand{\skScheitelwinkel}{\skizze{\draw[mass] (143.1:0.6) arc[start angle=143.1, end angle=201.8, radius=0.6];
	\draw[mass] (323.1:0.6) arc[start angle=323.1, end angle=381.8, radius=0.6];
	\draw[fig] (-2,-0.8) -- (2,0.8); \draw[fig] (-1.6,1.2) -- (1.6,-1.2);
	\node at (172:1.05) {$\alpha$}; \node at (352:1.1) {$\alpha'$};}}
\newcommand{\skStufenwinkel}{\skizze{\draw[mass] (0.5,0) arc[start angle=0, end angle=53.13, radius=0.5];
	\draw[mass] (1.7,1.6) arc[start angle=0, end angle=53.13, radius=0.5];
	\draw[fig] (-1.2,0) -- (2.6,0); \draw[fig] (-1.2,1.6) -- (2.6,1.6); \draw[fig] (-0.6,-0.8) -- (1.8,2.4);
	\node at (26:0.85) {$\alpha$}; \node at ($(1.2,1.6)+(26:0.85)$) {$\beta$};}}
\newcommand{\skWechselwinkel}{\skizze{\draw[mass] (0.5,0) arc[start angle=0, end angle=53.13, radius=0.5];
	\draw[mass] (0.7,1.6) arc[start angle=180, end angle=233.13, radius=0.5];
	\draw[fig] (-1.2,0) -- (2.6,0); \draw[fig] (-1.2,1.6) -- (2.6,1.6); \draw[fig] (-0.6,-0.8) -- (1.8,2.4);
	\node at (26:0.85) {$\alpha$}; \node at ($(1.2,1.6)+(206:0.85)$) {$\beta$};}}

% Abbildungen
\newcommand{\skSymmetrieachse}{\skizze{\draw[fig] (0,0) -- (2.6,0) -- (1.3,2.1) -- cycle;
	\draw[mass, hilf] (1.3,-0.5) -- (1.3,2.5);}}
\newcommand{\skSpiegelung}{\skizze{\draw[mass, hilf] (0,-0.4) -- (0,1.8);
	\draw[fig] (-2,0) -- (-0.5,0.3) -- (-1.5,1.4) -- cycle;
	\draw[fig, blue!45] (2,0) -- (0.5,0.3) -- (1.5,1.4) -- cycle;}}
\newcommand{\skVerschiebung}{\skizze{\draw[fig] (0,0) -- (1.3,0) -- (0.3,1.1) -- cycle;
	\draw[fig, blue!45] (2,0.6) -- (3.3,0.6) -- (2.3,1.7) -- cycle;
	\draw[mass, -{Stealth}] (0.3,1.1) -- (2.3,1.7);}}
\newcommand{\skDrehung}{\skizze{\draw[fig] (0.6,0) -- (1.9,0) -- (0.8,0.9) -- cycle;
	\draw[fig, blue!45] (0,0.6) -- (0,1.9) -- (-0.9,0.8) -- cycle;
	\fill (0,0) circle (2pt) node[below left] {$Z$};
	\draw[mass, -{Stealth}] (10:2.2) arc[start angle=10, end angle=80, radius=2.2];}}

% Dreiecke
\newcommand{\skDreieck}{\skizze{\draw[fig] (0,0) -- (3,0) -- (1,2) -- cycle;
	\node[below left] at (0,0) {$A$}; \node[below right] at (3,0) {$B$}; \node[above] at (1,2) {$C$};}}
\newcommand{\skRechtwinklig}{\skizze{\draw[mass] (0.5,0) arc[start angle=0, end angle=90, radius=0.5];
	\fill[red] (45:0.27) circle (1pt); \draw[fig] (0,0) -- (3,0) -- (0,2) -- cycle;}}
\newcommand{\skGleichschenklig}{\skizze{\draw[fig] (0,0) -- (2.6,0) -- (1.3,2.2) -- cycle;
	\draw[mass] (0.52,1.176) -- (0.78,1.024); \draw[mass] (1.82,1.024) -- (2.08,1.176);}}
\newcommand{\skGleichseitig}{\skizze{\draw[fig] (0,0) -- (2.6,0) -- (1.3,2.252) -- cycle;
	\draw[mass] (1.3,-0.15) -- (1.3,0.15); \draw[mass] (0.52,1.201) -- (0.78,1.051);
	\draw[mass] (1.82,1.051) -- (2.08,1.201);}}
\newcommand{\skHypotenuse}{\skizze{\draw[fig] (0,2) -- (0,0) -- (3,0); \draw[red, very thick] (3,0) -- (0,2);
	\draw[mass] (0.5,0) arc[start angle=0, end angle=90, radius=0.5]; \fill[red] (45:0.27) circle (1pt);
	\node[above right] at (1.5,1) {$c$};}}
\newcommand{\skKathete}{\skizze{\draw[fig] (3,0) -- (0,2); \draw[red, very thick] (0,2) -- (0,0) -- (3,0);
	\draw[mass] (0.5,0) arc[start angle=0, end angle=90, radius=0.5]; \fill[red] (45:0.27) circle (1pt);
	\node[left] at (0,1) {$b$}; \node[below] at (1.5,0) {$a$};}}
\newcommand{\skPythagoras}{\skizze{\draw[fig] (0,0) -- (2,0) -- (0,1.5) -- cycle;
	\draw[fig, fill=blue!10] (0,0) rectangle (2,-2); \draw[fig, fill=blue!10] (0,0) rectangle (-1.5,1.5);
	\draw[fig, fill=red!10] (2,0) -- (3.5,2) -- (1.5,3.5) -- (0,1.5) -- cycle;
	\node at (1,-1) {$a^2$}; \node at (-0.75,0.75) {$b^2$}; \node at (1.75,1.75) {$c^2$};}}
\newcommand{\skHoehe}{\skizze{\draw[fig] (0,0) -- (3.2,0) -- (1,2) -- cycle;
	\draw[mass, hilf] (1,2) -- node[right] {$h$} (1,0);
	\draw[mass] (1.35,0) arc[start angle=0, end angle=90, radius=0.35]; \fill[red] (1,0) ++(45:0.19) circle (0.8pt);}}

% Vierecke, Vielecke
\newcommand{\skViereck}{\skizze{\draw[fig] (0,0) -- (2.8,0.3) -- (2.3,1.9) -- (0.4,1.5) -- cycle;}}
\newcommand{\skQuadrat}{\skizze{\draw[fig] (0,0) rectangle (1.9,1.9);}}
\newcommand{\skRechteck}{\skizze{\draw[fig] (0,0) rectangle (3.2,1.6);}}
\newcommand{\skParallelogramm}{\skizze{\draw[fig] (0,0) -- (2.6,0) -- (3.4,1.5) -- (0.8,1.5) -- cycle;}}
\newcommand{\skRaute}{\skizze{\draw[fig] (0,0) -- (1.4,-1) -- (2.8,0) -- (1.4,1) -- cycle;}}
\newcommand{\skTrapez}{\skizze{\draw[fig] (0,0) -- (3.2,0) -- (2.4,1.5) -- (0.6,1.5) -- cycle;}}
\newcommand{\skDrachen}{\skizze{\draw[fig] (1,0) -- (2,1.6) -- (1,2.3) -- (0,1.6) -- cycle;
	\draw[mass, hilf] (1,-0.3) -- (1,2.6);}}
\newcommand{\skFuenfeck}{\skizze{\draw[fig] (90:1.2) -- (162:1.2) -- (234:1.2) -- (306:1.2) -- (18:1.2) -- cycle;}}
\newcommand{\skSechseck}{\skizze{\draw[fig] (0:1.2) -- (60:1.2) -- (120:1.2) -- (180:1.2) -- (240:1.2) -- (300:1.2) -- cycle;}}
\newcommand{\skFlaeche}{\skizze{\filldraw[gray, opacity=0.5] (0,0) rectangle (3,1.6);
	\draw[fig] (0,0) rectangle (3,1.6); \node at (1.5,0.8) {$A$};}}
\newcommand{\skUmfangViereck}{\skizze{\draw[red, very thick] (0,0) rectangle (3,1.6); \node at (1.5,0.8) {$u$};}}

% Kreis
\newcommand{\skKreis}{\skizze{\draw[fig] (0,0) circle (1.2); \fill (0,0) circle (1.5pt);}}
\newcommand{\skMittelpunkt}{\skizze{\draw[fig] (0,0) circle (1.2); \fill[red] (0,0) circle (2.5pt) node[below] {$M$};}}
\newcommand{\skRadius}{\skizze{\draw[fig] (0,0) circle (1.2); \fill (0,0) circle (1.5pt) node[below] {$M$};
	\draw[mass] (0,0) -- node[above] {$r$} (30:1.2);}}
\newcommand{\skDurchmesser}{\skizze{\draw[fig] (0,0) circle (1.2); \fill (0,0) circle (1.5pt);
	\draw[mass] (200:1.2) -- (20:1.2); \node[above] at (100:0.25) {$d$};}}
\newcommand{\skKreisumfang}{\skizze{\draw[red, very thick] (0,0) circle (1.2); \fill (0,0) circle (1.5pt);}}
\newcommand{\skSektor}{\skizze{\filldraw[gray, opacity=0.5] (0,0) -- (1.2,0) arc[start angle=0, end angle=70, radius=1.2] -- cycle;
	\draw[fig] (0,0) circle (1.2); \draw[mass] (1.2,0) -- (0,0) -- (70:1.2);}}
\newcommand{\skSehne}{\skizze{\draw[fig] (0,0) circle (1.2); \draw[mass] (160:1.2) -- (40:1.2);}}
\newcommand{\skTangente}{\skizze{\draw[fig] (0,0) circle (1.1); \draw[mass] ($(50:1.1)+(140:1.5)$) -- ($(50:1.1)+(-40:1.5)$);
	\fill[red] (50:1.1) circle (2.5pt);}}

% Körper (Schrägbild: Tiefe unter 45° und halb so lang, verdeckte Kanten gestrichelt)
\newcommand{\skWuerfel}{\skizze{\def\d{0.57}
	\draw[fig, hilf, thick] (0,0) -- (\d,\d) -- (1.6+\d,\d) (\d,\d) -- (\d,1.6+\d);
	\draw[fig] (0,0) rectangle (1.6,1.6); \draw[fig] (1.6,0) -- (1.6+\d,\d) -- (1.6+\d,1.6+\d) -- (\d,1.6+\d) -- (0,1.6)
		(1.6,1.6) -- (1.6+\d,1.6+\d);}}
\newcommand{\skQuader}{\skizze{\def\d{0.6}
	\draw[fig, hilf, thick] (0,0) -- (\d,\d) -- (2.6+\d,\d) (\d,\d) -- (\d,1.3+\d);
	\draw[fig] (0,0) rectangle (2.6,1.3); \draw[fig] (2.6,0) -- (2.6+\d,\d) -- (2.6+\d,1.3+\d) -- (\d,1.3+\d) -- (0,1.3)
		(2.6,1.3) -- (2.6+\d,1.3+\d);}}
\newcommand{\skPrisma}{\skizze{
	\draw[fig, hilf, thick] (0,0) -- (1.4,0.7) -- (3,0.7) (1.4,0.7) -- (2.2,2);
	\draw[fig] (0,0) -- (1.6,0) -- (0.8,1.3) -- cycle;
	\draw[fig] (1.6,0) -- (3,0.7) -- (2.2,2) -- (0.8,1.3);}}
\newcommand{\skZylinder}{\skizze{
	\draw[fig, hilf, thick] (1,0) arc[start angle=0, end angle=180, x radius=1, y radius=0.3];
	\draw[fig] (-1,0) arc[start angle=180, end angle=360, x radius=1, y radius=0.3];
	\draw[fig] (0,1.9) ellipse[x radius=1, y radius=0.3];
	\draw[fig] (-1,0) -- (-1,1.9) (1,0) -- (1,1.9);}}
\newcommand{\skKegel}{\skizze{
	\draw[fig, hilf, thick] (1,0) arc[start angle=0, end angle=180, x radius=1, y radius=0.3];
	\draw[fig] (-1,0) arc[start angle=180, end angle=360, x radius=1, y radius=0.3];
	\draw[fig] (-1,0) -- (0,2.2) -- (1,0);}}
\newcommand{\skPyramide}{\skizze{
	\draw[fig, hilf, thick] (0,0) -- (0.8,0.6) -- (2.8,0.6) (0.8,0.6) -- (1.4,2.3);
	\draw[fig] (0,0) -- (2,0) -- (2.8,0.6) -- (1.4,2.3) -- cycle (2,0) -- (1.4,2.3);}}
\newcommand{\skKugel}{\skizze{\draw[fig] (0,0) circle (1.1);
	\draw[fig, hilf, thick] (1.1,0) arc[start angle=0, end angle=180, x radius=1.1, y radius=0.3];
	\draw[fig] (-1.1,0) arc[start angle=180, end angle=360, x radius=1.1, y radius=0.3];}}
\newcommand{\skKante}{\skizze{\def\d{0.57}
	\draw[fig, hilf, thick] (0,0) -- (\d,\d) -- (1.6+\d,\d) (\d,\d) -- (\d,1.6+\d);
	\draw[fig] (0,0) rectangle (1.6,1.6); \draw[fig] (1.6,0) -- (1.6+\d,\d) -- (1.6+\d,1.6+\d) -- (\d,1.6+\d) -- (0,1.6);
	\draw[red, very thick] (1.6,1.6) -- (1.6+\d,1.6+\d);}}
\newcommand{\skEcke}{\skizze{\def\d{0.57}
	\draw[fig, hilf, thick] (0,0) -- (\d,\d) -- (1.6+\d,\d) (\d,\d) -- (\d,1.6+\d);
	\draw[fig] (0,0) rectangle (1.6,1.6); \draw[fig] (1.6,0) -- (1.6+\d,\d) -- (1.6+\d,1.6+\d) -- (\d,1.6+\d) -- (0,1.6)
		(1.6,1.6) -- (1.6+\d,1.6+\d); \fill[red] (1.6+\d,1.6+\d) circle (2.5pt);}}
\newcommand{\skSeitenflaeche}{\skizze{\def\d{0.57}
	\draw[fig, hilf, thick] (0,0) -- (\d,\d) -- (1.6+\d,\d) (\d,\d) -- (\d,1.6+\d);
	\fill[gray, opacity=0.5] (0,0) rectangle (1.6,1.6);
	\draw[fig] (0,0) rectangle (1.6,1.6); \draw[fig] (1.6,0) -- (1.6+\d,\d) -- (1.6+\d,1.6+\d) -- (\d,1.6+\d) -- (0,1.6)
		(1.6,1.6) -- (1.6+\d,1.6+\d);}}
\newcommand{\skNetz}{\skizze{\begin{scope}[scale=0.62]
	\foreach \x/\y in {1/0, 0/1, 1/1, 2/1, 3/1, 1/2} {\draw[fig, fill=blue!8] (\x,\y) rectangle ++(1,1);}
	\end{scope}}}

% Funktionen (Achsen dünn mit Pfeil, Graph blau, Hervorhebungen rot)
\newcommand{\achsen}[4]{\draw[achse] (#1,0) -- (#2,0) node[below left] {$x$}; \draw[achse] (0,#3) -- (0,#4) node[below left] {$y$};}
\newcommand{\skKoordinaten}{\skizze{\draw[black!25, very thin] (-1.5,-1.5) grid (2.5,2); \achsen{-1.6}{2.7}{-1.6}{2.2}
	\fill[red] (1,1.5) circle (2pt) node[right] {$P$};}}
\newcommand{\skGraph}{\skizze{\achsen{-1.8}{2.4}{-0.6}{2.4}
	\begin{scope}\clip (-1.8,-0.6) rectangle (2.2,2.2); \draw[fig, domain=-1.8:2.2, samples=40] plot (\x, {0.3*\x*\x*\x - 0.6*\x + 0.9});\end{scope}}}
\newcommand{\skXAchse}{\skizze{\draw[achse, red, thick] (-1.8,0) -- (2.4,0) node[below left] {$x$}; \draw[achse] (0,-1.2) -- (0,1.6) node[below left] {$y$};}}
\newcommand{\skParabel}{\skizze{\achsen{-1.9}{1.9}{-0.6}{2.4}
	\begin{scope}\clip (-1.9,-0.6) rectangle (1.9,2.2); \draw[fig, domain=-1.8:1.8, samples=40] plot (\x, {0.7*\x*\x});\end{scope}}}
\newcommand{\skScheitelpunkt}{\skizze{\achsen{-0.6}{3}{-0.6}{2.4}
	\begin{scope}\clip (-0.6,-0.6) rectangle (3,2.2); \draw[fig, domain=-0.3:3, samples=40] plot (\x, {0.7*(\x-1.4)*(\x-1.4)+0.4});\end{scope}
	\fill[red] (1.4,0.4) circle (2.5pt) node[below] {$S$};}}
\newcommand{\skNullstelle}{\skizze{\achsen{-1.9}{1.9}{-1.2}{1.8}
	\begin{scope}\clip (-1.9,-1.2) rectangle (1.9,1.6); \draw[fig, domain=-1.8:1.8, samples=40] plot (\x, {0.6*\x*\x-0.9});\end{scope}
	\fill[red] (-1.2247,0) circle (2.5pt) (1.2247,0) circle (2.5pt);}}
\newcommand{\skLinear}{\skizze{\achsen{-0.8}{3}{-0.6}{2.4}
	\begin{scope}\clip (-0.8,-0.6) rectangle (3,2.2); \draw[fig, domain=-0.8:3] plot (\x, {0.5*\x+0.5});\end{scope}}}
\newcommand{\skSteigung}{\skizze{\achsen{-0.8}{3}{-0.6}{2.4}
	\begin{scope}\clip (-0.8,-0.6) rectangle (3,2.2); \draw[fig, domain=-0.8:3] plot (\x, {0.5*\x+0.5});\end{scope}
	\draw[mass] (0.6,0.8) -- node[below] {$\Delta x$} (2.2,0.8) -- node[right] {$\Delta y$} (2.2,1.6);}}
\newcommand{\skAchsenabschnitt}{\skizze{\achsen{-0.8}{3}{-0.6}{2.4}
	\begin{scope}\clip (-0.8,-0.6) rectangle (3,2.2); \draw[fig, domain=-0.8:3] plot (\x, {0.5*\x+0.5});\end{scope}
	\fill[red] (0,0.5) circle (2.5pt) node[left] {$n$};}}
\newcommand{\skSchnittpunkt}{\skizze{\draw[fig] (-0.2,0) -- (3,1.6); \draw[fig] (-0.2,1.8) -- (3,0.2);
	\fill[red] (1.6,0.9) circle (2.5pt) node[above=2pt] {$S$};}}
\newcommand{\skExponential}{\skizze{\achsen{-1.9}{1.9}{-0.6}{2.4}
	\begin{scope}\clip (-1.9,-0.6) rectangle (1.9,2.2); \draw[fig, domain=-1.9:1.9, samples=40] plot (\x, {0.5*2^\x});\end{scope}}}
\newcommand{\skLogarithmus}{\skizze{\achsen{-0.6}{3.2}{-1.4}{1.6}
	\begin{scope}\clip (-0.6,-1.4) rectangle (3.2,1.4); \draw[fig, domain=0.2:3.2, samples=40] plot (\x, {ln(\x)});\end{scope}}}
\newcommand{\sinusgraph}{\achsen{-0.3}{3.5}{-1.2}{1.4}\draw[fig, domain=0:3.3, samples=50] plot (\x, {0.9*sin(2*\x r)});}
\newcommand{\skSinus}{\skizze{\sinusgraph}}
\newcommand{\skKosinus}{\skizze{\achsen{-0.3}{3.5}{-1.2}{1.4}\draw[fig, domain=0:3.3, samples=50] plot (\x, {0.9*cos(2*\x r)});}}
\newcommand{\skHochpunkt}{\skizze{\sinusgraph \fill[red] (0.785,0.9) circle (2.5pt) node[above] {$H$};}}
\newcommand{\skTiefpunkt}{\skizze{\sinusgraph \fill[red] (2.356,-0.9) circle (2.5pt) node[below] {$T$};}}
\newcommand{\skWendepunkt}{\skizze{\sinusgraph \fill[red] (1.571,0) circle (2.5pt) node[above right] {$W$};}}
\newcommand{\skMonotonSteigend}{\skizze{\achsen{-0.4}{3}{-0.4}{2.2}\draw[fig, domain=0:2.7, samples=40] plot (\x, {0.25*\x*\x+0.2});}}
\newcommand{\skMonotonFallend}{\skizze{\achsen{-0.4}{3}{-0.4}{2.2}\draw[fig, domain=0.25:2.8, samples=40] plot (\x, {0.5/\x+0.1});}}
\newcommand{\skAchsensymmetrisch}{\skizze{\draw[mass, hilf] (0,-0.6) -- (0,2.2); \achsen{-1.9}{1.9}{-0.6}{2.4}
	\begin{scope}\clip (-1.9,-0.6) rectangle (1.9,2.2); \draw[fig, domain=-1.8:1.8, samples=40] plot (\x, {0.3*\x*\x*\x*\x - 0.9*\x*\x + 1});\end{scope}}}
\newcommand{\skPunktsymmetrisch}{\skizze{\achsen{-1.9}{1.9}{-1.4}{1.6}
	\begin{scope}\clip (-1.9,-1.4) rectangle (1.9,1.4); \draw[fig, domain=-1.8:1.8, samples=40] plot (\x, {0.4*\x*\x*\x - 0.9*\x});\end{scope}
	\fill[red] (0,0) circle (2.5pt);}}
\newcommand{\skAsymptote}{\skizze{\achsen{-0.4}{3.2}{-0.4}{2.4}\draw[mass, hilf] (-0.3,0.4) -- (3.1,0.4);
	\draw[fig, domain=0.3:3.1, samples=40] plot (\x, {0.5/\x+0.4});}}
\newcommand{\skIntegral}{\skizze{\fill[gray, opacity=0.5] (0.5,0) -- plot[domain=0.5:2.5, samples=30] (\x, {0.3*(\x-1.5)*(\x-1.5)+0.6}) -- (2.5,0) -- cycle;
	\achsen{-0.4}{3.2}{-0.4}{2.2}\draw[fig, domain=0.1:2.9, samples=40] plot (\x, {0.3*(\x-1.5)*(\x-1.5)+0.6});
	\node[below] at (0.5,0) {$a$}; \node[below] at (2.5,0) {$b$};}}
\newcommand{\skTangenteGraph}{\skizze{\achsen{-0.4}{3.2}{-0.4}{2.4}
	\begin{scope}\clip (-0.4,-0.4) rectangle (3.2,2.2); \draw[fig, domain=0:3, samples=40] plot (\x, {0.25*\x*\x});
	\draw[mass] (0.4,-0.4) -- (3.2,2.4);\end{scope} \fill[red] (2,1) circle (2pt);}}
\newcommand{\skFolge}{\skizze{\achsen{-0.3}{3.4}{-0.3}{2.2}
	\foreach \n in {1,...,10} {\fill[blue] ({0.3*\n},{1.6-1.2/\n}) circle (1.6pt);}
	\draw[mass, hilf] (0,1.6) -- (3.2,1.6);}}

% Vektoren
\newcommand{\skVektor}{\skizze{\draw[mass, very thick, -{Stealth}] (0,0) -- node[above left] {$\vec{a}$} (2.6,1.3);}}
\newcommand{\skOrtsvektor}{\skizze{\achsen{-0.4}{3}{-0.4}{2.2}\draw[mass, very thick, -{Stealth}] (0,0) -- (2.2,1.5);
	\fill (2.2,1.5) circle (2pt) node[above] {$P$}; \node[below left] at (0,0) {$O$};}}
\newcommand{\skKreuzprodukt}{\skizze{\draw[fig, -{Stealth}] (0,0) -- (2.2,0) node[below] {$\vec{a}$};
	\draw[fig, -{Stealth}] (0,0) -- (1,-0.8) node[below] {$\vec{b}$};
	\draw[mass, very thick, -{Stealth}] (0,0) -- (0,2) node[right] {$\vec{a}\times\vec{b}$};}}
\newcommand{\skEbene}{\skizze{\filldraw[blue!10, draw=blue, very thick] (0,0) -- (2.8,0) -- (3.6,1.4) -- (0.8,1.4) -- cycle;
	\node at (1.8,0.7) {$E$};}}

% Mengen (wie im Skill latex-venn-diagramme: Schnitt über \clip)
\newcommand{\venn}{\draw[fig] (0,0) circle (1) (1.2,0) circle (1); \node at (-0.5,0) {$A$}; \node at (1.7,0) {$B$};}
\newcommand{\skSchnittmenge}{\skizze{\begin{scope}\clip (0,0) circle (1); \fill[gray, opacity=0.5] (1.2,0) circle (1);\end{scope}\venn}}
\newcommand{\skVereinigung}{\skizze{\fill[gray, opacity=0.5] (0,0) circle (1) (1.2,0) circle (1);\venn}}
\newcommand{\skTeilmenge}{\skizze{\draw[fig] (0,0) ellipse[x radius=1.6, y radius=1.1]; \draw[fig, fill=gray!40] (-0.5,-0.1) circle (0.55);
	\node at (-0.5,-0.1) {$A$}; \node at (0.8,0.3) {$B$};}}

% Stochastik
\newcommand{\skBaum}{\skizze{\draw[fig, thick] (0,0) -- (1.4,0.9) (0,0) -- (1.4,-0.9) (1.4,0.9) -- (2.8,1.35) (1.4,0.9) -- (2.8,0.45)
	(1.4,-0.9) -- (2.8,-0.45) (1.4,-0.9) -- (2.8,-1.35);
	\foreach \p in {(0,0),(1.4,0.9),(1.4,-0.9),(2.8,1.35),(2.8,0.45),(2.8,-0.45),(2.8,-1.35)} {\fill[blue] \p circle (1.6pt);}}}
\newcommand{\skVierfeldertafel}{{\scriptsize\begin{tabular}{c|cc|c}
	& $B$ & $\overline{B}$ & \\ \hline $A$ & & & \\ $\overline{A}$ & & & \\ \hline & & & $1$\end{tabular}}}
\newcommand{\skWertetabelle}{{\scriptsize\begin{tabular}{c|c|c|c|c} $x$ & $0$ & $1$ & $2$ & $3$ \\ \hline $y$ & $1$ & $3$ & $5$ & $7$\end{tabular}}}
\newcommand{\skSaeulen}{\skizze{\draw[achse] (-0.2,0) -- (3.4,0); \draw[achse] (0,-0.2) -- (0,2.2);
	\foreach \k/\h in {1/0.4, 2/1.2, 3/1.8, 4/1.2, 5/0.4} {\filldraw[fill=blue!25, draw=blue, thick] ({0.55*\k-0.2},0) rectangle ++(0.4,\h);}}}

% Ober- und Untersumme (Ergänzung aus Teil II)
\newcommand{\rechtecksumme}[2]{\skizze{\foreach \i in {0,...,3} {
		\pgfmathsetmacro{\xl}{0.4+0.6*\i} \pgfmathsetmacro{\xr}{1.0+0.6*\i}
		\pgfmathsetmacro{\h}{0.18*(#1)*(#1)+0.3}
		\filldraw[fill=#2!30, draw=#2] (\xl,0) rectangle (\xr,\h);}
	\achsen{-0.3}{3.2}{-0.3}{2.4}\draw[fig, domain=0:2.9, samples=30] plot (\x, {0.18*\x*\x+0.3});}}
\newcommand{\skUntersumme}{\rechtecksumme{\xl}{gray}}
\newcommand{\skObersumme}{\rechtecksumme{\xr}{red}}

%----------------------------------------------------------------------------------------------------%
% Ergänzungen: Geometrie, Körper, analytische Geometrie
%----------------------------------------------------------------------------------------------------%
\newcommand{\skScheitel}{\skizze{\draw[fig] (2.3,0) -- (0,0) -- (50:2.3); \fill[red] (0,0) circle (3pt) node[below] {$S$};}}
\newcommand{\skSchenkel}{\skizze{\draw[fig] (0,0) -- (50:2.3); \draw[red, very thick] (0,0) -- (2.3,0); \fill (0,0) circle (1.5pt);}}
\newcommand{\skKomplement}{\skizze{\draw[mass] (0.8,0) arc[start angle=0, end angle=35, radius=0.8];
	\draw[mass] (35:1.1) arc[start angle=35, end angle=90, radius=1.1];
	\draw[fig] (2.3,0) -- (0,0) -- (0,2.3); \draw[fig] (0,0) -- (35:2.3);
	\node at (17:1.25) {$\alpha$}; \node at (62:1.5) {$\beta$};}}
\newcommand{\skEntgegen}{\skizze{\draw[mass] (0.5,0) arc[start angle=0, end angle=53.13, radius=0.5];
	\draw[mass] ($(1.2,1.6)+(233.13:0.45)$) arc[start angle=233.13, end angle=360, radius=0.45];
	\draw[fig] (-1.2,0) -- (2.6,0); \draw[fig] (-1.2,1.6) -- (2.6,1.6); \draw[fig] (-0.6,-0.8) -- (1.8,2.4);
	\node at (26:0.85) {$\alpha$}; \node at ($(1.2,1.6)+(300:0.8)$) {$\beta$};}}
\newcommand{\skMittelsenkrechte}{\skizze{\draw[fig] (0,0) -- (3,0); \fill (0,0) circle (2pt) (3,0) circle (2pt);
	\draw[mass] (1.5,-0.8) -- (1.5,1.4); \draw[mass] (1.85,0) arc[start angle=0, end angle=90, radius=0.35];
	\fill[red] (1.5,0) ++(45:0.19) circle (0.8pt); \draw[mass] (0.75,-0.13) -- (0.75,0.13) (2.25,-0.13) -- (2.25,0.13);}}
\newcommand{\skWinkelhalbierende}{\skizze{\draw[mass] (0.8,0) arc[start angle=0, end angle=30, radius=0.8];
	\draw[mass] (30:0.95) arc[start angle=30, end angle=60, radius=0.95];
	\draw[fig] (2.3,0) -- (0,0) -- (60:2.3); \draw[mass, hilf] (0,0) -- (30:2.6);}}
\newcommand{\skInnenwinkel}{\skizze{\draw[mass] (0.55,0) arc[start angle=0, end angle=63.43, radius=0.55];
	\draw[mass] ($(3,0)+(135:0.55)$) arc[start angle=135, end angle=180, radius=0.55];
	\draw[mass] ($(1,2)+(243.43:0.5)$) arc[start angle=243.43, end angle=315, radius=0.5];
	\draw[fig] (0,0) -- (3,0) -- (1,2) -- cycle;}}
\newcommand{\skAussenwinkel}{\skizze{\draw[mass] (3.6,0) arc[start angle=0, end angle=135, radius=0.6];
	\draw[fig] (0,0) -- (3,0) -- (1,2) -- cycle; \draw[fig, hilf, thick] (3,0) -- (4.3,0);
	\node at ($(3,0)+(60:0.98)$) {$\delta$};}}
\newcommand{\skBasis}{\skizze{\draw[fig] (0,0) -- (1.3,2.2) -- (2.6,0); \draw[red, very thick] (0,0) -- (2.6,0);}}
\newcommand{\skBasiswinkel}{\skizze{\draw[mass] (0.6,0) arc[start angle=0, end angle=59.4, radius=0.6];
	\draw[mass] ($(2.6,0)+(120.6:0.6)$) arc[start angle=120.6, end angle=180, radius=0.6];
	\draw[fig] (0,0) -- (2.6,0) -- (1.3,2.2) -- cycle;}}
\newcommand{\rwdreieck}{\draw[mass] (0.85,0) arc[start angle=0, end angle=30.96, radius=0.85]; \node at (14:1.2) {$\alpha$};
	\draw[black!60, thick] (3,0.4) arc[start angle=90, end angle=180, radius=0.4]; \fill[black!60] ($(3,0)+(135:0.2)$) circle (0.8pt);}
\newcommand{\skAnkathete}{\skizze{\rwdreieck \draw[fig] (0,0) -- (3,1.8) -- (3,0); \draw[red, very thick] (0,0) -- (3,0);}}
\newcommand{\skGegenkathete}{\skizze{\rwdreieck \draw[fig] (3,0) -- (0,0) -- (3,1.8); \draw[red, very thick] (3,0) -- (3,1.8);}}
\newcommand{\skSeitenhalbierende}{\skizze{\draw[fig] (0,0) -- (3,0) -- (1,2) -- cycle; \draw[mass] (1,2) -- (1.5,0);
	\draw[mass] (0.75,-0.13) -- (0.75,0.13) (2.25,-0.13) -- (2.25,0.13);}}
\newcommand{\skInkreis}{\skizze{\draw[fig] (0,0) -- (3,0) -- (1,2) -- cycle; \draw[mass] (1.204,0.744) circle (0.744);}}
\newcommand{\skUmkreis}{\skizze{\draw[mass] (1.5,0.5) circle (1.581); \draw[fig] (0,0) -- (3,0) -- (1,2) -- cycle;}}
\newcommand{\skDiagonale}{\skizze{\draw[fig] (0,0) rectangle (3,1.6); \draw[mass] (0,0) -- (3,1.6);}}
\newcommand{\skNeck}[1]{\skizze{\draw[fig] (90:1.2) \foreach \k in {1,...,#1} { -- ({90+\k*360/#1}:1.2) } -- cycle;}}
\newcommand{\skSekante}{\skizze{\draw[fig] (0,0) circle (1.2); \draw[mass] ($(150:1.2)!-0.35!(30:1.2)$) -- ($(150:1.2)!1.35!(30:1.2)$);
	\fill[red] (150:1.2) circle (1.6pt) (30:1.2) circle (1.6pt);}}
\newcommand{\skBeruehrpunkt}{\skizze{\draw[fig] (0,0) circle (1.1); \draw[fig, thick] ($(50:1.1)+(140:1.5)$) -- ($(50:1.1)+(-40:1.5)$);
	\fill[red] (50:1.1) circle (3pt) node[above right] {$B$};}}
\newcommand{\skKreisbogen}{\skizze{\draw[fig] (0,0) circle (1.2); \draw[red, line width=1.8pt] (20:1.2) arc[start angle=20, end angle=110, radius=1.2];}}
\newcommand{\skKreisabschnitt}{\skizze{\fill[gray, opacity=0.5] (150:1.2) arc[start angle=150, end angle=30, radius=1.2] -- cycle;
	\draw[fig] (0,0) circle (1.2); \draw[mass] (150:1.2) -- (30:1.2);}}
\newcommand{\skKreisflaeche}{\skizze{\filldraw[fill=gray!50, draw=blue, very thick] (0,0) circle (1.2);}}

% Körper
\newcommand{\wuerfelbasis}{\def\d{0.57}\draw[fig, hilf, thick] (0,0) -- (\d,\d) -- (1.6+\d,\d) (\d,\d) -- (\d,1.6+\d);
	\draw[fig] (0,0) rectangle (1.6,1.6); \draw[fig] (1.6,0) -- (1.6+\d,\d) -- (1.6+\d,1.6+\d) -- (\d,1.6+\d) -- (0,1.6) (1.6,1.6) -- (1.6+\d,1.6+\d);}
\newcommand{\quaderbasis}{\def\d{0.6}\draw[fig, hilf, thick] (0,0) -- (\d,\d) -- (2.6+\d,\d) (\d,\d) -- (\d,1.3+\d);
	\draw[fig] (0,0) rectangle (2.6,1.3); \draw[fig] (2.6,0) -- (2.6+\d,\d) -- (2.6+\d,1.3+\d) -- (\d,1.3+\d) -- (0,1.3) (2.6,1.3) -- (2.6+\d,1.3+\d);}
\newcommand{\zylinderbasis}{\draw[fig, hilf, thick] (1,0) arc[start angle=0, end angle=180, x radius=1, y radius=0.3];
	\draw[fig] (-1,0) arc[start angle=180, end angle=360, x radius=1, y radius=0.3];
	\draw[fig] (0,1.9) ellipse[x radius=1, y radius=0.3]; \draw[fig] (-1,0) -- (-1,1.9) (1,0) -- (1,1.9);}
\newcommand{\skFlaechendiagonale}{\skizze{\wuerfelbasis \draw[red, very thick] (0,0) -- (1.6,1.6);}}
\newcommand{\skRaumdiagonale}{\skizze{\wuerfelbasis \draw[red, very thick] (0,0) -- (1.6+\d,1.6+\d);}}
\newcommand{\skGrundflaeche}{\skizze{\quaderbasis \fill[red, opacity=0.4] (0,0) -- (2.6,0) -- (2.6+\d,\d) -- (\d,\d) -- cycle;}}
\newcommand{\skDeckflaeche}{\skizze{\quaderbasis \fill[red, opacity=0.4] (0,1.3) -- (2.6,1.3) -- (2.6+\d,1.3+\d) -- (\d,1.3+\d) -- cycle;}}
\newcommand{\skMantel}{\skizze{\fill[red, opacity=0.35] (-1,0) arc[start angle=180, end angle=360, x radius=1, y radius=0.3] -- (1,1.9)
	arc[start angle=0, end angle=-180, x radius=1, y radius=0.3] -- cycle; \zylinderbasis}}
\newcommand{\skPyramidenstumpf}{\skizze{\draw[fig, hilf, thick] (0,0) -- (0.8,0.6) -- (2.8,0.6) (0.8,0.6) -- (1.1,1.45);
	\draw[fig] (0,0) -- (2,0) -- (2.8,0.6) -- (2.1,1.45) -- (1.1,1.45) -- (0.7,1.15) -- cycle (0.7,1.15) -- (1.7,1.15) -- (2.1,1.45) (1.7,1.15) -- (2,0);}}
\newcommand{\skKegelstumpf}{\skizze{\draw[fig, hilf, thick] (1.1,0) arc[start angle=0, end angle=180, x radius=1.1, y radius=0.3];
	\draw[fig] (-1.1,0) arc[start angle=180, end angle=360, x radius=1.1, y radius=0.3];
	\draw[fig] (0,1.5) ellipse[x radius=0.5, y radius=0.15]; \draw[fig] (-1.1,0) -- (-0.5,1.5) (1.1,0) -- (0.5,1.5);}}
\newcommand{\skTetraeder}{\skizze{\draw[fig, hilf, thick] (0,0) -- (2.6,0.3);
	\draw[fig] (0,0) -- (1.6,-0.4) -- (2.6,0.3) -- (1.3,1.9) -- cycle (1.6,-0.4) -- (1.3,1.9);}}
\newcommand{\skHohlzylinder}{\skizze{\zylinderbasis \draw[fig] (0,1.9) ellipse[x radius=0.5, y radius=0.15];
	\draw[fig, hilf, thin] (-0.5,0) -- (-0.5,1.9) (0.5,0) -- (0.5,1.9);}}
\newcommand{\skSchieferZylinder}{\skizze{\draw[fig, hilf, thick] (1,0) arc[start angle=0, end angle=180, x radius=1, y radius=0.3];
	\draw[fig] (-1,0) arc[start angle=180, end angle=360, x radius=1, y radius=0.3];
	\draw[fig] (0.8,1.8) ellipse[x radius=1, y radius=0.3]; \draw[fig] (-1,0) -- (-0.2,1.8) (1,0) -- (1.8,1.8);}}
\newcommand{\drehachse}[2]{\draw[dash dot] (0,#1) -- (0,#2); \draw[mass, -{Stealth[length=3pt]}] (-0.3,#2-0.15) arc[start angle=180, end angle=500, x radius=0.3, y radius=0.1];}
\newcommand{\skRotation}{\skizze{\filldraw[fill=gray!50, draw=blue, very thick] (0,-1.1) arc[start angle=-90, end angle=90, radius=1.1] -- cycle;
	\draw[fig, hilf, thick] (0,1.1) arc[start angle=90, end angle=270, radius=1.1]; \drehachse{-1.4}{1.6}}}
\newcommand{\skKugelsegment}{\skizze{\fill[gray, opacity=0.5] (27.04:1.1) arc[start angle=27.04, end angle=152.96, radius=1.1] -- cycle;
	\draw[fig] (0,0) circle (1.1); \draw[fig] (-0.98,0.5) arc[start angle=180, end angle=360, x radius=0.98, y radius=0.2];
	\draw[fig, hilf, thick] (0.98,0.5) arc[start angle=0, end angle=180, x radius=0.98, y radius=0.2];}}
\newcommand{\skKugelsektor}{\skizze{\fill[gray, opacity=0.5] (0,0) -- (55:1.1) arc[start angle=55, end angle=125, radius=1.1] -- cycle;
	\draw[fig] (0,0) circle (1.1); \draw[fig, thick] (55:1.1) -- (0,0) -- (125:1.1);
	\draw[fig, thick] (-0.63,0.9) arc[start angle=180, end angle=360, x radius=0.63, y radius=0.13];}}
\newcommand{\skEllipsoid}{\skizze{\draw[fig] (0,0) ellipse[x radius=1.5, y radius=0.9];
	\draw[fig, hilf, thick] (1.5,0) arc[start angle=0, end angle=180, x radius=1.5, y radius=0.3];
	\draw[fig] (-1.5,0) arc[start angle=180, end angle=360, x radius=1.5, y radius=0.3];}}
\newcommand{\skParaboloid}{\skizze{\draw[fig, domain=-1.3:1.3, samples=30] plot (\x, {0.9*\x*\x});
	\draw[fig] (0,1.521) ellipse[x radius=1.3, y radius=0.22];}}
\newcommand{\skHyperboloid}{\skizze{\draw[fig, domain=-1.2:1.2, samples=30, variable=\t] plot ({0.55*cosh(\t)}, {0.9*sinh(\t)});
	\draw[fig, domain=-1.2:1.2, samples=30, variable=\t] plot ({-0.55*cosh(\t)}, {0.9*sinh(\t)});
	\draw[fig] (0,1.358) ellipse[x radius=0.996, y radius=0.2];
	\draw[fig] (-0.996,-1.358) arc[start angle=180, end angle=360, x radius=0.996, y radius=0.2];
	\draw[fig, hilf, thick] (0.996,-1.358) arc[start angle=0, end angle=180, x radius=0.996, y radius=0.2];}}

% Trigonometrie und analytische Geometrie
\newcommand{\skEinheitskreis}{\skizze{\achsen{-1.6}{1.9}{-1.5}{1.8}\draw[fig] (0,0) circle (1.1); \node[below right] at (1.05,0) {$1$};}}
\newcommand{\skBogenmass}{\skizze{\draw[fig] (1.6,0) -- (0,0) -- (60:1.6); \draw[red, very thick] (1.6,0) arc[start angle=0, end angle=60, radius=1.6];
	\node at (30:2.0) {$b$}; \node[below] at (0.8,0) {$r$}; \node at (30:0.6) {$\alpha$};}}
\newcommand{\skPolar}{\skizze{\achsen{-0.4}{3}{-0.4}{2.4}\draw[mass] (0.9,0) arc[start angle=0, end angle=36, radius=0.9];
	\draw[mass, very thick] (0,0) -- node[above left, inner sep=1pt] {$r$} (2.2,1.6); \fill (2.2,1.6) circle (2pt) node[above] {$P$};
	\node at (17:1.35) {$\varphi$};}}
\newcommand{\skKreisgleichung}{\skizze{\achsen{-1.7}{1.9}{-1.6}{1.8}\draw[fig] (0,0) circle (1.2); \draw[mass] (0,0) -- node[above left, inner sep=1pt] {$r$} (35:1.2);}}
\newcommand{\ellipsebasis}{\achsen{-2}{2.2}{-1.4}{1.6}\draw[fig] (0,0) ellipse[x radius=1.6, y radius=1];}
\newcommand{\skEllipse}{\skizze{\ellipsebasis}}
\newcommand{\skBrennpunkt}{\skizze{\ellipsebasis \fill[red] (-1.249,0) circle (2.5pt) node[below] {$F_1$} (1.249,0) circle (2.5pt) node[below] {$F_2$};}}
\newcommand{\skHalbachse}{\skizze{\ellipsebasis \draw[red, very thick] (0,0) -- node[below] {$a$} (1.6,0) (0,0) -- node[left] {$b$} (0,1);}}
\newcommand{\skHyperbel}{\skizze{\achsen{-2.2}{2.4}{-1.6}{1.8}
	\draw[fig, domain=-1.6:1.6, samples=30, variable=\t] plot ({0.7*cosh(\t)}, {0.6*sinh(\t)});
	\draw[fig, domain=-1.6:1.6, samples=30, variable=\t] plot ({-0.7*cosh(\t)}, {0.6*sinh(\t)});}}
\newcommand{\skLeitlinie}{\skizze{\achsen{-1.2}{2.4}{-1.7}{1.9}
	\draw[fig, domain=-1.5:1.5, samples=30, variable=\t] plot ({\t*\t/1.6}, {\t});
	\draw[mass, hilf] (-0.4,-1.6) -- (-0.4,1.7) node[left] {$l$}; \fill (0.4,0) circle (2pt) node[below right] {$F$};}}
